% BHU PYQ -- Manually transcribed & error-corrected
\documentclass[11pt,a4paper]{article}

\usepackage[a4paper,top=2.5cm,bottom=2.5cm,left=2.8cm,right=2.8cm,headheight=14pt]{geometry}
\usepackage{amsmath,amssymb}
\usepackage[T1]{fontenc}
\usepackage[utf8]{inputenc}
\usepackage{lmodern}
\usepackage{microtype}
\usepackage{fancyhdr}
\usepackage{enumitem}
\usepackage{hyperref}
\usepackage{lastpage}
\usepackage{parskip}

\hypersetup{colorlinks=true,urlcolor=black,linkcolor=black,
  pdftitle={CHB-503: Organic Chemistry-III -- BHU PYQ},pdfauthor={Banaras Hindu University}}

\pagestyle{fancy}\fancyhf{}
\renewcommand{\headrulewidth}{0.4pt}\renewcommand{\footrulewidth}{0.4pt}
\fancyhead[L]{\small\textit{Banaras Hindu University}}
\fancyhead[R]{\small\textit{CHB-503: Organic Chemistry-III}}
\fancyfoot[C]{\small Page \thepage\ of \pageref{LastPage}}

\newlist{parts}{enumerate}{2}
\setlist[parts,1]{label=(\alph*),leftmargin=*,itemsep=5pt,topsep=3pt}
\setlist[parts,2]{label=(\roman*),leftmargin=*,itemsep=3pt,topsep=2pt}
\newcommand{\pts}[1]{\hfill\textit{[#1]}}

\begin{document}

\begin{center}
  {\large\bfseries BANARAS HINDU UNIVERSITY}\\[2pt]
  {\normalsize B.Sc. (Hons.) Semester V Examination, 2016-17}\\[6pt]
  \rule{0.8\linewidth}{0.5pt}\\[6pt]
  {\large\bfseries Paper No. CHB-503: Organic Chemistry-III}\\[4pt]
  {\normalsize Time: 3 Hours\hfill Maximum Marks: 70}
\end{center}
\rule{\linewidth}{0.4pt}
\small
\noindent\textit{Instructions: Answer five questions in total, including Question~1 which is compulsory.}
\normalsize
\bigskip

\noindent\textbf{Question 1.} Answer all parts of the following: \pts{5 $\times$ 2 = 10}
\begin{parts}
  \item Predict the product: Nitrobenzene $\xrightarrow{\text{Na/liq. NH}_3, \text{EtOH}} ?$
  \item Describe radical halogenation of alkanes with a suitable example.
  \item How can cyclopropanes be prepared using carbenes?
  \item Give a preparation method for cyclobutane from malonic acid.
  \item What is the Wittig reaction? Give one example.
\end{parts}

\medskip\rule{\linewidth}{0.2pt}

\noindent\textbf{Question 2.}
\begin{parts}
  \item Explain the Benzidine rearrangement reaction. Provide evidences to support that it is an intramolecular rearrangement reaction. \pts{7}
  \item Discuss the conditions and mechanisms to obtain simple and modified Wittig reaction products. \pts{7}
\end{parts}

\medskip\rule{\linewidth}{0.2pt}

\noindent\textbf{Question 3.}
\begin{parts}
  \item Explain the concept of a plane of symmetry and center of symmetry with suitable chiral and achiral examples. \pts{5}
  \item Define resolution of racemic modification and explain two chemical methods of resolution. \pts{5}
  \item Explain regioselectivity with a suitable example. \pts{4}
\end{parts}

\medskip\rule{\linewidth}{0.2pt}

\noindent\textbf{Question 4.}
\begin{parts}
  \item What do you understand by asymmetric induction? Explain Cram's rule of asymmetric induction using a suitable carbonyl reaction. \pts{6}
  \item Explain prochirality with suitable examples. How can trigonal carbons be made into stereogenic centers? Explain with a reaction. \pts{6}
  \item Write the mechanism of the Pinacol-pinacolone rearrangement. \pts{2}
\end{parts}

\vfill
\rule{\linewidth}{0.4pt}
\begin{center}\small
  \href{https://bhu-exam.vercel.app/}{Click here to visit our website}\\[4pt]
  \href{https://me-aryan.vercel.app/}{Click here to visit our contributor}\\[4pt]
  \href{https://quashan-me.vercel.app/}{Click here to visit our contributor}
\end{center}

\end{document}
